% ============ §26 Separability (原书 pp.198-207) ============
\section{Separability}\label{sec:26}\index{separable algebra, section 26}\index{separably generated extension, section 26}\index{linearly disjoint, section 26}\index{p-basis, section 26}\index{differential basis, section 26}\index{imperfection module, section 26}\index{Cartier equality, section 26}\index{0-étale, section 26}

Let $k$ be a field and $A$ a $k$-algebra. We say that $A$ is \emph{separable} over $k$ if for
every extension field $k'$ of $k$, the ring $A' = A \otimes_k k'$ is reduced, that is does
not contain nilpotents. From the definition, one sees at once the following:
\begin{enumerate}[(1)]
  \item a subalgebra of a separable $k$-algebra is separable;
  \item $A$ is separable over $k$ if and only if every finitely generated $k$-
        subalgebra of $A$ is separable over $k$;
  \item for $A$ to be separable over $k$ it is sufficient that $A \otimes_k k'$ is reduced
        for every finitely generated extension field $k'$ of $k$;
  \item if $A$ is separable over $k$ and $k'$ is an extension field of $k$ then
        $A \otimes_k k'$ is separable over $k'$.
\end{enumerate}

\begin{remark}
When $A$ is a finite $k$-algebra, the separability condition can be checked using
the discriminant. The \emph{trace} of an element $a$ of $A$, denoted by $\tr_{A/k}(a)$, is
defined to be the trace of the $k$-linear mapping $A \to A$ induced by
multiplication by $a$. Let $\omega_1, \ldots, \omega_n$ be a linear basis of $A$ over $k$. Then
$d = \det(\tr_{A/k}(\omega_i\omega_j))$ is called a \emph{discriminant} of $A$ over $k$. If we use
another basis $\omega'_1, \ldots, \omega'_n$, and if $\omega_i = \sum_j c_{ij}\omega_j$, then the
discriminant with respect to this basis is $\det(c_{ij})^2 \cdot d$. Thus $d = 0$ or
$d \ne 0$ is a property of $A$ independent of the choice of basis. Now we claim that
$A$ is separable if and only if $d \ne 0$.

If $k'$ is an extension field of $k$ and $A' = A \otimes_k k'$, then $\omega_1, \ldots, \omega_n$
is also a linear basis of $A'$ over $k'$, and so $d$ is also a discriminant of $A'$
over $k'$. If $A'$ is not reduced, let $N = \nil(A')$. Take a basis $\omega'_1, \ldots, \omega'_n$
of $A'$ such that $\omega'_1, \ldots, \omega'_r$ span $N$. Then $\omega'_i\omega'_j$ is nilpotent for
$i \le r$, hence its trace is zero. It follows that $\det(\tr(\omega'_i\omega'_j)) = 0$, and so
$d = 0$. Conversely, if $A$ is separable over $k$, take an algebraic closure $K$ of $k$.
Then $A \otimes_k K$ is reduced. Therefore we need only prove that if $k$ is
algebraically closed and $A$ is reduced then $d \ne 0$. Now $A$ is an Artinian reduced
ring, hence is a finite product of fields, each of which (as a finite extension of
$k$) is isomorphic to $k$. Thus $A = ke_1 + \cdots + ke_n$ with $e_ie_j = 0$ for $i \neq j$ and
$e_i^2 = e_i$. Hence $\tr(e_i) = 1$ and $d \ne 0$.
\end{remark}

In what follows we consider mainly the case when $A$ is a field. If $K$ is an
algebraic extension field of $k$, and is separable in the usual sense (that is
every element of $K$ is a root of a polynomial with coefficients in $k$ having no
multiple roots), then $K$ is separable over $k$ in our sense. To see this, by (2)
above we can assume that $K$ is finitely generated over $k$, and then according to
the well-known primitive element theorem of field theory, $K \cong k[X]/(f(X))$,
where $f \in k[X]$ is irreducible and with no multiple roots. Then if $k'$ is an
extension field of $k$ we have
\[
  K \otimes_k k' = k'[X]/(f(X)),
\]
and when we factorise $f$ into primes in $k'[X]$ we get $f = f_1 \cdots f_s$ with
$(f_i, f_j) = 1$ for $i \neq j$, so that by Theorem~\ref{thm:1.4},
\[
  k'[X]/(f) = k'[X]/(f_1) \times \cdots \times k'[X]/(f_s);
\]
since this is a direct product of fields, it is reduced.

We say that an extension field $K$ of $k$ is \emph{separably generated} over $k$ if $K$ has
a separating transcendence basis over $k$, that is a transcendence basis $T$ such
that $K$ is a separable algebraic extension of $k(T)$.

\begin{theorem}\label{thm:26.1}
A separably generated extension field is separable.
\end{theorem}

\begin{proof}
Let $k$ be a field and $K$ a separably generated extension of $k$, with $T$ a
separating transcendence basis of $K$. If $k'$ is any extension field of $k$ then
$k(T) \otimes_k k'$ is a ring of fractions of $k[T] \otimes_k k' = k'[T]$, so that it is an
integral domain with field of fractions $k'(T)$. Thus
$K \otimes_k k' = K \otimes_{k(T)} (k(T) \otimes_k k')$ is a subring of $K \otimes_k k'(T)$. Now
$K$ is a separable algebraic extension of $k(T)$, so that as we have seen above
$K \otimes_{k(T)} k'(T)$ is a direct product of fields, hence reduced; therefore so is
$K \otimes_k k'$.
\end{proof}

\begin{theorem}\label{thm:26.2}
Let $k$ be a field of characteristic $p$, and $K$ a finitely generated extension
field of $k$; then the following conditions are equivalent:
\begin{enumerate}[(1)]
  \item $K$ is separable over $k$;
  \item $K \otimes_k k^{1/p}$ is reduced;
  \item $K$ is separably generated over $k$.
\end{enumerate}
\end{theorem}

\begin{proof}
(1) $\Leftarrow$ (3) has just been proved, and (3) $\Rightarrow$ (2) is clear from the
proof of Theorem~\hyperref[thm:26.1]{1}, since $k^{1/p}$ is an extension field of $k$.

(1) $\Rightarrow$ (3) Let $K = k(x_1, \ldots, x_n)$; we can assume that $x_1, \ldots, x_r$ is a
transcendence basis for $K$ over $k$. Assume furthermore that $x_{r+1}, \ldots, x_q$ are
separable algebraic over $k(x_1, \ldots, x_r)$, and that $x_{q+1}$ is not; set $y = x_{q+1}$
and let $f(Y^p)$ be the minimal polynomial of $y$ over $k(x_1, \ldots, x_r)$. The
coefficients of $f(Y^p)$ are rational functions of $x_1, \ldots, x_r$, so that clearing
denominators we get an irreducible polynomial
$F(X_1, \ldots, X_r, Y^p) \in k[X_1, \ldots, X_r, Y]$, with $F(x, y^p) = 0$. Now if
$\partial F/\partial X_i = 0$ for $1 \le i \le r$ then $F(X, Y)$ is the $p$-th power of a polynomial
$G(X, Y)$ with coefficients in $k^{1/p}$, but then we would have
\[
  k[x_1, \ldots, x_r, y] \otimes_k k^{1/p} = (k[X, Y]/(F(X, Y))) \otimes_k k^{1/p}
      = k^{1/p}[X, Y]/(G(X, Y)^p);
\]
this is a subring of $K \otimes_k k^{1/p}$ containing nilpotent elements, and this
contradicts (2). Hence we can assume that $\partial F/\partial X_1 \ne 0$. Then $x_1$ is
separable algebraic over $k(x_2, \ldots, x_r, y)$, and hence so are $x_{r+1}, \ldots, x_q$.
Therefore exchanging $x_1$ and $y = x_{q+1}$, we find that $x_{r+1}, \ldots, x_{q+1}$ are
separably algebraic over $k(x_1, \ldots, x_r)$, so that by induction on $q$, we have
(3).
\end{proof}

\begin{remark}
As we have seen in the proof, if $K = k(x_1, \ldots, x_n)$ is separable over $k$ then we
can choose a separating transcendence basis from among $x_1, \ldots, x_n$.
\end{remark}

\begin{theorem}\label{thm:26.3}
If $k$ is a perfect field then every extension field $K$ of $k$ is separable over
$k$, and a $k$-algebra $A$ is separable if and only if it is reduced.
\end{theorem}

\begin{proof}
Recall that a field $k$ is perfect if every algebraic extension of $k$ is separable.
If $k$ has characteristic $0$ then every extension field $K$ is separably generated,
and therefore separable. In characteristic $p$, perfect implies $k = k^{1/p}$, so that
by the previous theorem, all subfields of $K$ finitely generated over $k$ are
separable, so that $K$ itself is separable over $k$. (Caution: $K$ may fail to be
separably generated over $k$; for a counter-example, let $x$ be an indeterminate
over $k$, and $K = k(x, x^{p^{-1}}, x^{p^{-2}}, \ldots)$.) Now we show that if $A$ is a
reduced $k$-algebra then $A$ is separable. We can assume that $A$ is finitely
generated over $k$. Then $A$ is a Noetherian ring, and the total ring of fractions
$K$ of $A$ is of the form $K = K_1 \times \cdots \times K_r$ by Ex.~\ref{ex:6.5-ex}. Each $K_i$ is
separable over $k$, so that $K$ is also separable, and hence so is its subring $A$.
\end{proof}

In general two subfields $K$, $K'$ of a given field $L$ are said to be \emph{linearly
disjoint}\index{linearly disjoint} over a common subfield $k$ if the following three equivalent
conditions are satisfied:
\begin{enumerate}[(a)]
  \item if $\alpha_1, \ldots, \alpha_n \in K$ are linearly independent over $k$ they are also
        linearly independent over $K'$;
  \item the same with $K$ and $K'$ interchanged;
  \item if we write $K[K']$ for the subring of $L$ generated by $K$ and $K'$, the
        natural map $K \otimes_k K' \to K[K']$ is an isomorphism.
\end{enumerate}

\textit{Proof of equivalence.} (a) $\Rightarrow$ (c) Let $\xi = \sum x_i \otimes y_i$ be an element of
the kernel of $K \otimes_k K' \to K[K']$. Suppose that $x_1, \ldots, x_r$ are linearly
independent over $k$, and that the remainder $x_{r+1}, \ldots, x_n$ are linear
combinations of them, and rewrite $\xi = \sum x_i \otimes y_i'$. The image in $K[K']$ of
$\xi$ is $\sum x_iy_i'$, but if this is $0$ then by (a) we have $y_i' = 0$ for all $i$, so that
$\xi = 0$. This proves (c).

(c) $\Rightarrow$ (a) is also easy; finally, since (c) is symmetric in $K$ and $K'$ we of
course also get (a) $\Leftrightarrow$ (b).

Let $k$ be a field of characteristic $p$, and $K$ an extension field of $k$. Inside
the algebraic closure of $K$, consider the subfields $k^{p^{-n}} =
\{\alpha \in \bar{K} \mid \alpha^{p^n} \in k\}$; these are purely inseparable extension fields of
$k$, and $\bigcup_n k^{p^{-n}}$ is the smallest perfect field containing $k$.

\begin{theorem}[S.~MacLane]\label{thm:26.4}
Let $k$ and $K$ be as above. We have
\begin{enumerate}[(i)]
  \item if $K$ is separable over $k$ then $K$ and $k^{p^{-n}}$ are linearly disjoint over $k$;
  \item if $K$ and $k^{p^{-n}}$ are linearly disjoint over $k$ for some $n > 0$ then $K$ is
        separable over $k$.
\end{enumerate}
\end{theorem}

\begin{proof}
(i) It is enough to show that $K$ and $k^{p^{-n}}$ are linearly disjoint for every $n$.
Suppose that $\alpha_1, \ldots, \alpha_m \in k^{p^{-n}}$ are linearly independent over $k$. If
$\sum \xi_i\alpha_i = 0$ with $\xi_i \in K$, then we have to show that $\xi_i = 0$ for all $i$. Since
$k$ and $K$ are as in the statement, we can assume that $\xi_1, \ldots, \xi_m$ are
linearly independent over $k$. If we set $k_1 = k(\alpha_1, \ldots, \alpha_m)$, so that $k_1$ is a
finite extension of $k$; for some sufficiently large $n$ we have $k^{p^{-n}} \supset k_1$, and
if we set $A = K \otimes_k k_1$, then $A$ is a reduced ring. However, $A$ is finite as a
$K$-module, so is a zero-dimensional ring, but the $p^n$-th power of any element of
$A$ is in $K$, so that we see that $A$ has only one prime ideal. Hence $A$ is a field,
and is isomorphic to $K[k_1]$. From this we get $\sum \xi_i \otimes \alpha_i = 0$, that is
$\xi_i = 0$ for all $i$.

(ii) If $K$ and $k^{p^{-n}}$ are linearly disjoint, then $k^{p^{-n}} \subset k^{p^{-m}}$, and hence
$K$ and $k^{p^{-1}}$ are also linearly disjoint over $k$, so that $K \otimes_k k^{p^{-1}}$ is a
field. If $K'$ is a subfield of $K$ which is finitely generated over $k$ then by
Theorem~\hyperref[thm:26.2]{2}, $K'$ is separable over $k$. Hence $K$ is also separable over $k$.
\end{proof}

Let $K$ be an extension field of a field $k$; then $\Omega_{K/k}$ is a vector space over
$K$, and is generated by $\{dx \mid x \in K\}$, so that there exists a subset $B \subset K$
such that $\{dx \mid x \in B\}$ forms a basis of the vector space $\Omega_{K/k}$. A subset
$B \subset K$ with this property is called a \emph{differential basis} of $K$ over $k$. The
following condition (*) is necessary and sufficient for a subset
$\{x_\lambda\}_{\lambda \in \Lambda} \subset K$ to form a differential basis for $K$ over $k$.

\textbf{(*)} if $y_\lambda \in K$ are specified for every $\lambda \in \Lambda$ in an arbitrary way, then
there exists a unique $D \in \Der_k(K)$ such that $D(x_\lambda) = y_\lambda$ for all $\lambda$.

For $x_1, \ldots, x_n \in K$, let us study the condition for $dx_1, \ldots, dx_n \in \Omega_{K/k}$ to
be linearly independent over $K$. If $k$ has characteristic $0$ then this is equivalent
to $x_1, \ldots, x_n$ being algebraically independent over $k$. Indeed, if there exists
$f(X) \in k[X_1, \ldots, X_n]$ such that $f(x_1, \ldots, x_n) = 0$, then choose such a
relation of smallest degree; suppose for instance that $X_1$ actually appears in $f$,
so that $f_1 = \partial f/\partial X_1$ is non-zero, but of smaller degree than $f$, and hence
$f_1(x) \ne 0$. Then $f(x) = 0$ gives
\[
  0 = df = \sum_i f_i(x)dx_i,
\]
so that $dx_1, \ldots, dx_n$ are linearly dependent. Conversely, if $x_1, \ldots, x_n$ are
algebraically independent over $k$ then there exists a transcendence basis $B$ of
$K/k$ containing these, so that there exists $k$-derivations $D_i$ of the purely
transcendental extension $k(B)$ satisfying
\[
  D_i(x_i) = 1 \quad \text{and} \quad D_i(x) = 0 \text{ for } x \in B, x \neq x_i,
\]
(namely, $\partial/\partial x_i$). Moreover, $K$ is a separable algebraic extension of $k(B)$, and
so by Theorem~\hyperref[thm:25.3]{25.3} is 0-\'etale, so that the derivations $D_i$ extend to
derivations from $K$ to $K$. Then since $D_i(x_j) = \delta_{ij}$, the differentials
$dx_1, \ldots, dx_n \in \Omega_{K/k}$ are clearly linearly independent. Thus in this case a
differential basis is the same thing as a transcendence basis.

Now consider a field $k$ of characteristic $p$. We say that elements
$x_1, \ldots, x_n \in K$ of an extension field $K$ are \emph{$p$-independent} over $k$ if
$[K^p(k, x_1, \ldots, x_n) : K^p(k)] = p^n$, and a subset $B \subset K$ is $p$-independent if any
finite subset of $B$ is $p$-independent. This condition means precisely that the
set
\[
  \Gamma_B = \left\{ x_1^{\alpha_1} \cdots x_n^{\alpha_n} \;\middle|\;
      \begin{array}{l}
        x_1, \ldots, x_n \text{ are distinct} \\
        \text{elements of } B \text{ and } 0 \le \alpha_i < p
      \end{array} \right\}
\]
is linearly independent over $K^p(k)$; the elements of $\Gamma_B$ are called the
\emph{$p$-monomials} of $B$. If $B$ is not $p$-independent, we say it is $p$-dependent.
The condition of $p$-independence is not just a property of $B$ and $k$, but also
depends on $K$. If $B \subset K$ is $p$-independent over $k$ and $K = K^p(k, B)$, we say
that $B$ is a \emph{$p$-basis} of $K/k$. If $C \subset K$ is $p$-independent over $k$ then one
can easily show by Zorn's lemma that there exists a $p$-basis of $K/k$ containing
$C$.

One sees easily that $B$ is a $p$-basis of $K/k$ is equivalent to $\Gamma_B$ being a
basis of $K$ over $K^p(k)$ in the sense of linear algebra. If this holds then any map
$D\colon B \to K$ has a unique extension to an element $D \in \Der_k(K)$. Indeed, for a
$p$-monomial of $B$ we set
\[
  D(x_1^{\alpha_1} \cdots x_n^{\alpha_n}) = \sum_{r=1}^n \alpha_r x_1^{\alpha_1} \cdots
      x_r^{\alpha_r - 1} \cdots x_n^{\alpha_n}D(x_r),
\]
and extend $D$ to $K$ as a $K^p(k)$-linear map; then $D$ is a $k$-derivation. Thus a
$p$-basis $B$ is a differential basis of $K/k$. Conversely, if $B'$ is a differential
basis of $K/k$ then $B'$ is $p$-independent over $k$; for if $x_1, \ldots, x_n \in B'$ are
$p$-dependent, we can assume that $x_1 \in K^p(k, x_2, \ldots, x_n)$, so that we can
write $x_1 = f(x_2, \ldots, x_n)$, where $f$ is a polynomial with coefficients in $K^p(k)$.
Then in $\Omega_{K/k}$ we get $dx_1 = \sum_{i \ge 2}(\partial f/\partial x_i)dx_i$, which
contradicts the linear independence of $dx_1, \ldots, dx_n$. Now if we take a $p$-basis
$B$ of $K/k$ containing $B'$, then since both $B$ and $B'$ are differential bases, we
have $B = B'$. We summarise the above as follows:

\begin{theorem}\label{thm:26.5}
The notion of differential basis coincides with transcendence basis in
characteristic $0$, and with $p$-basis in characteristic $p$.
\end{theorem}

Now we look at the relation between separability and differential bases. Letting
$\pi \subset k$ denote the prime subfield of $k$, we write $\Omega_k$ for $\Omega_{k/\pi}$.

\begin{theorem}\label{thm:26.6}
For a field extension $K/k$, the following conditions are equivalent:
\begin{enumerate}[(1)]
  \item $K/k$ is separable;
  \item for any subfield $k' \subset k$ the standard map $\Omega_{k/k'} \otimes_k K \to \Omega_{K/k'}$
        is injective;
  \item[$2'$] for any subfield $k' \subset k$ and any differential basis $B$ of $k/k'$, there
        exists a differential basis of $K/k'$ containing $B$;
  \item $\Omega_k \otimes_k K \to \Omega_K$ is injective;
  \item any derivation of $k$ to an arbitrary $K$-module $M$ extends to a
        derivation from $K$ to $M$.
\end{enumerate}
\end{theorem}

\begin{proof}
(2) and (2$'$) are clearly equivalent, and (2) $\Rightarrow$ (3) $\Rightarrow$ (4) are trivial. In
characteristic $0$, both (1) and (2$'$) hold, so that we need only consider the case
$p > 0$.

(1) $\Rightarrow$ (2$'$) Since $K$ and $k^{1/p}$ are linearly disjoint over $k$, we can apply the
isomorphism $x \mapsto x^p$ to all three of these to get $K^p$ and $k$ linearly disjoint
over $k^p$. Hence $K^p(k^p, k') = K^p(k')$ and $k$ are linearly disjoint over
$k^p(k')$ (see Ex.~\ref{ex:26.1-ex} below). If we choose a $p$-basis $B$ of $k$ over $k'$ then
the set $\Gamma_B$ of $p$-monomials of $B$ is linearly independent over $k^p(k')$, hence
also linearly independent over $K^p(k')$, and $B$ as a subset of $K$ is also
$p$-independent over $k'$. Therefore $B$ can be extended to a $p$-basis of $K/k'$.

(2$'$) $\Rightarrow$ (1) If we take a $p$-basis $B$ of $k$ over $\pi$ then the set $\Gamma_B$ of
$p$-monomials is a basis of $k$ over $k^p$. $\{dx \mid x \in B\}$ is a basis of $\Omega_k$ over
$k$, and by assumption is linearly independent in $\Omega_K$ over $K$, so that
$\Gamma_B$ is also linearly independent over $K^p$. Therefore
\[
  k \otimes_{k^p} K^p \cong k(K^p),
\]
and $k$ and $K^p$ are linearly disjoint over $k^p$, so that by
Theorem~\hyperref[thm:26.4]{4}, $K/k$ is separable.
\end{proof}

Let $k$ be a field of characteristic $p$, and $\pi \subset k$ the prime subfield; a $p$-basis
of $k$ over $\pi$ is called an \emph{absolute $p$-basis} of $k$. If $k_1 \subset k$ is any perfect field
contained in $k$ then $k^p(k_1) = k^p = k^p(\pi)$, so that an absolute $p$-basis of $k$ is
the same as a $p$-basis over any perfect subfield of $k$.

\begin{theorem}\label{thm:26.7}
Let $k$ be a field of characteristic $p$, and $K$ an extension field of $k$. If an
absolute $p$-basis of $k$ is also an absolute $p$-basis of $K$, then $K$ is 0-\'etale over
$k$, and conversely.
\end{theorem}

\begin{proof}
Consider a commutative diagram of ring homomorphisms
\[
\begin{tikzcd}
K \arrow[r, "\bar{u}"] & \bar{C} = C/N \\
k \arrow[u, "\subset"] \arrow[ru, "j"'] & C \arrow[u]
\end{tikzcd}
\]
where $C = C/N$ is as in the definition with $N$ an ideal of $C$ satisfying $N^2 = 0$,
and $g$ the natural map. For $\alpha \in K$, if we choose $a \in C$ such that
$\bar{u}(\alpha) = g(a)$, then $a^p$ is independent of the choice of $a$. For if
$g(a) = g(a')$ then we can write $a' = a + x$ with $x \in N$, and since $C$ is a $k$-algebra
via $j$ it is a ring of characteristic $p$, so that
\[
  a'^p = a^p + x^p = a^p.
\]
Now we define a map $v_0\colon K^p \to C$ by $v_0(\alpha^p) = a^p$ for $\alpha \in K$, where $a$ is
such that $\bar{u}(\alpha) = g(a)$; one checks easily that $v_0$ is a homomorphism, and
coincides with $j$ on $k^p$. So far we have not used the assumption on $K/k$. Now
since by assumption $K$ is separable over $k$, and $K = K^p[k]$, we can think of $K$ as
\[
  K = K^p \otimes_{k^p} k,
\]
and thus we can define $v\colon K \to C$ by letting $v$ be equal to $v_0$ on $K^p$, and
equal to $j$ on $k$; this is a lifting of $\bar{u}$ to $C$. Uniqueness of the lifting is
clear from the fact that $K^p(k) = K$, so that $\Omega_{K/k} = 0$.

Conversely, if we assume that $K/k$ is 0-\'etale, then first of all, from
0-unramified we have $\Omega_{K/k} = 0$, so that by 0-smoothness and
Theorem~\hyperref[thm:25.1]{25.1} we have $\Omega_K = \Omega_k \otimes_k K$. Thus an absolute $p$-basis of $k$ is
also an absolute $p$-basis of $K$.
\end{proof}

\begin{theorem}\label{thm:26.8}
Let $K/k$ be a separable extension of fields of characteristic $p$, and let $B$ be a
$p$-basis of $K/k$. Then $B$ is algebraically independent over $k$, and $K$ is 0-\'etale
over $k(B)$.
\end{theorem}

\begin{proof}
Suppose by contradiction that $b_1, \ldots, b_n \in B$ are algebraically dependent
over $k$. Suppose that $0 \ne f \in k[X_1, \ldots, X_n]$ is a polynomial of minimal degree
among all those with $f(b) = 0$, and set $\deg f = d$. Then write
\[
  f(X_1, \ldots, X_n) = \sum_{0 \le i_\nu < p} g_{i_1, \ldots, i_n}(X_1^p, \ldots, X_n^p)
      X_1^{i_1} \cdots X_n^{i_n}.
\]
Then since $b_1, \ldots, b_n$ are $p$-independent over $k$ and $f(b) = 0$, we have
$g_{i_1, \ldots, i_n}(b^p) = 0$ for all values of $i_1, \ldots, i_n$. However, since
\[
  d \ge \deg g_{i_1, \ldots, i_n}(X^p) + i_1 + \cdots + i_n,
\]
by choice of $f$ we must have $f(X) = g_{0, \ldots, 0}(X^p)$. Hence we can write $f$ in the
form $f(X) = h(X)^p$, with $h \in k^{1/p}[X_1, \ldots, X_n]$. However, since $K$ and $k^{1/p}$
are linearly disjoint over $k$, the monomials of degree $< d$ in $b_1, \ldots, b_n$, being
linearly independent over $k$, must also be linearly independent over $k^{1/p}$.
Thus $h(b) \ne 0$, but this contradicts $h(b)^p = f(b) = 0$. For the second half, see
the proof of the following theorem.
\end{proof}

\begin{remark}
Although $k(B)$ is purely transcendental over $k$, it does not follow that $K$ is
algebraic over $k(B)$; for a counter-example, let $K = k(x, x^{p^{-1}}, x^{p^{-2}}, \ldots)$,
with $x$ an indeterminate over $k$. In this case $B = \emptyset$.
\end{remark}

\begin{theorem}\label{thm:26.9}
If $K$ is a separable field extension of a field $k$, then $K$ is 0-smooth over $k$,
and conversely.
\end{theorem}

\begin{proof}
Let $B$ be a differential basis of $K/k$. If $K/k$ is separable then by
Theorems~\hyperref[thm:26.5]{5} and \hyperref[thm:26.8]{8}, $k(B)$ is purely transcendental over $k$. Then, as
one sees at once from the definition, $k(B)$ is 0-smooth over $k$. Moreover,
$K/k(B)$ is 0-\'etale. In characteristic $0$ this follows from
Theorem~\hyperref[thm:25.3]{25.3}; in characteristic $p$, by
Theorem~\hyperref[thm:26.6]{6} we have an exact sequence
$0 \to \Omega_k \otimes K \to \Omega_K \to \Omega_{K/k} \to 0$, so that putting together an absolute
$p$-basis of $k$ with $B$ we have an absolute $p$-basis of $K$, and this is clearly
also an absolute $p$-basis of $k(B)$. Hence $K/k(B)$ is 0-\'etale by
Theorem~\hyperref[thm:26.7]{7}. Therefore $K/k$ is 0-smooth.

Conversely, if $K/k$ is 0-smooth then by Theorem~\hyperref[thm:25.1]{25.1}, $\Omega_k \otimes K \to \Omega_K$
is injective, so that by Theorem~\hyperref[thm:26.6]{6}, $K/k$ is separable.
\end{proof}

\medskip
\noindent\textbf{Imperfection modules and the Cartier equality}

Quite generally, if $k \to A \to B$ are ring homomorphisms, we write $I_{B/A/k}$ for
the kernel of $\Omega_{A/k} \otimes_A B \to \Omega_{B/k}$, and call it the \emph{imperfection module} of
the $A$-algebra $B$ over $k$. If $k = \mathbb{Z}$ or $k = \mathbb{Z}/p$ (the prime field of characteristic
$p$) we omit $k$, and write $I_{B/A}$.

\begin{lemma}\label{lem:26.1}
Let $k \to K \to L \to L'$ be field homomorphisms. Then there exists a natural
exact sequence
\[
  0 \to I_{L/K/k} \otimes_K L' \to I_{L'/K/k} \to I_{L'/L/k}
      \to \Omega_{L/K} \otimes_L L' \to \Omega_{L'/K} \to \Omega_{L'/L} \to 0.
\]
\end{lemma}

\begin{proof}
We have a commutative diagram with exact rows
\[
\begin{tikzcd}
0 \arrow[r] & I_{L/K/k} \otimes_K L' \arrow[r] \arrow[d] &
\Omega_{K/k} \otimes_K L' \arrow[r] \arrow[d, "="] &
\Omega_{L/k} \otimes_L L' \arrow[r] \arrow[d, "\cong"] &
\Omega_{L/K} \otimes_L L' \arrow[r] \arrow[d] & 0 \\
0 \arrow[r] & I_{L'/K/k} \arrow[r] &
\Omega_{K/k} \otimes_K L' \arrow[r] &
\Omega_{L'/k} \arrow[r] &
\Omega_{L'/K} \arrow[r] & 0
\end{tikzcd}
\]
We abbreviate this as
\[
\begin{tikzcd}
0 \arrow[r] & X \arrow[r] \arrow[d, "i"] & A \arrow[r] \arrow[d, "="] & B \arrow[r] \arrow[d, "f_1"] & P \arrow[r] \arrow[d, "f_2"] & 0 \\
0 \arrow[r] & Y \arrow[r] & A \arrow[r] & C \arrow[r] & Q \arrow[r] & 0
\end{tikzcd}
\]
and from it construct
\[
\begin{tikzcd}
0 \arrow[r] & A/X \arrow[r] \arrow[d] & B \arrow[r] \arrow[d, "f_1"] & P \arrow[r] \arrow[d, "f_2"] & 0 \\
0 \arrow[r] & A/Y \arrow[r] & C \arrow[r] & Q \arrow[r] & 0
\end{tikzcd}
\]
applying the snake lemma gives the exact sequence
\[
  0 \to Y/X \to \Ker f_1 \to \Ker f_2 \to 0,
\]
from which we easily get
\[
  0 \to X \to Y \to \Ker f_1 \to P \to Q \to \Coker f_2 \to 0.
\]
This is just what we wanted to prove.
\end{proof}

\begin{theorem}[the Cartier equality]\label{thm:26.10}
Let $k$ be a perfect field, $K$ an extension of $k$ and $L$ a finitely generated
extension field of $K$; then
\[
  (*) \quad \dim_L \Omega_{L/K} = \trdeg_k L + \rank_L I_{L/K/k},
\]
(where $\rank_L$ denotes the dimension of a vector space over $L$).
\end{theorem}

\begin{proof}
Suppose that $k \subset K \subset L \subset L'$, with $L$ finitely generated over $K$ and $L'$
finitely generated over $L$. If the theorem holds for $k \subset L \subset L'$ and for
$k \subset K \subset L$ then both $I_{L'/L/k}$ and $I_{L/K/k} \otimes_K L$ are finite-dimensional over
$L$, so that by the lemma, $I_{L'/K/k}$ is also finite-dimensional, and
\begin{align*}
  \dim_{L'} \Omega_{L'/K} - \dim_{L'} \Omega_{L'/k}
      &= (\dim_{L'} \Omega_{L'/L} - \dim_{L'} \Omega_{L/k}) +
          (\dim_{L'} \Omega_{L/K} - \dim_L \Omega_{L/k}) \\
      &= \trdeg_k L' + \trdeg_K L = \trdeg_k L';
\end{align*}
thus the theorem also holds for $k \subset K \subset L'$. Now every finitely generated
field extension can be obtained by a succession of the following three kinds of
simple extensions:
\begin{enumerate}[(1)]
  \item $L = K(\alpha)$ where $\alpha$ is transcendental over $K$;
  \item $L = K(\alpha)$ where $\alpha$ is separable algebraic over $K$;
  \item $L = K(\alpha)$ where $\chark K = p$, and $\alpha^p = a \in K$, but $\alpha \notin K$.
\end{enumerate}
Hence we need only prove (*) in each of these special cases. (1) and (2) are
easy. For (3), if we write $L = K[X]/(X^p - a)$ we see that
\[
  \Omega_{L/k} = (\Omega_{K[X]/k} \otimes_{K[X]} L)/L\,da
      = (\Omega_{K/k} \otimes_K L \oplus L\,da)/L\,da,
\]
and $d\alpha \ne 0$. Furthermore, since $k$ is a perfect field, we have
$a \notin K^p = kK^p$, so that in $\Omega_{K/k}(= \Omega_K)$ we have $da \ne 0$,
$\rank \Omega_{L/K} = 1$ and $\rank I_{L/K/k} = 1$, so that (*) also holds in this case.
\end{proof}

\begin{remark}[\textbf{Harper's theorem}]
An ideal $I$ of a ring $R$ is called a \emph{differential ideal} if it maps into itself under
every derivation of $R$ to $R$. A ring $R$ is said to be \emph{differentiably simple} if it has
no non-trivial differential ideals. The following beautiful theorem is due to
L.~Harper, Jr.

\begin{theorem*}
A Noetherian ring $R$ of characteristic $p$ is differentiably simple if and only if
it has the form $R = k[T_1, \ldots, T_n]/(T_1^p, \ldots, T_n^p)$, where $k$ is a field of
characteristic $p$.
\end{theorem*}

The `if' part is easy. The proof of the `only if' part is not so easy and we refer
the reader to Harper [1] and Yuan [1]. Recently this theorem was used by
Kimura-Niitsuma [1] to prove the following theorem which had been known as
Kunz' Conjecture:

\begin{theorem*}
Let $R$ be a regular local ring of characteristic $p$, and let $S$ be a local subring
of $R$ containing $R^p$. Assume that $R$ is a finite $S$-module. Then $R$ has a
$p$-basis over $S$ if and only if $S$ is regular.
\end{theorem*}
\end{remark}

% ============ 习题 ============
\section*{Exercises to \S\ref{sec:26}}
\emph{Prove the following propositions.}

\begin{exercise}\label{ex:26.1-ex}
Let $L$ be a field and $k$, $k'$, $K$, $K'$ subfields of $L$; assume that $k \subset k' \subset K$
and $k \subset K'$, and that $K$ and $K'$ are linearly disjoint over $k$. Then we have
(i) $K \cap K' = k$, and (ii) $K$ and $k'(K')$ are linearly disjoint over $k'$.
\end{exercise}

\begin{exercise}\label{ex:26.2-ex}
Let $L$ be a separable extension of a field $K$; then $L((T_1, \ldots, T_n))$ is separable
over $K((T_1, \ldots, T_n))$. Here $L((T_1, \ldots, T_n))$ denotes the field of fractions of
$L[[T_1, \ldots, T_n]]$.
\end{exercise}
